\documentclass[11pt]{article}

% Pacotes
\RequirePackage{luatex85}
\usepackage{amsmath}
\usepackage{amssymb}
\usepackage{graphicx}
\usepackage[colorlinks=true, allcolors=blue]{hyperref}
\usepackage{amsthm}
\usepackage{framed, color,xcolor}
\usepackage{stmaryrd}
\usepackage{rotating}           % Usado para rotacionar o texto
\usepackage[all,knot,arc,import,poly]{xy}   % Pacote para desenhos gráficos
% Este pacote pode conflitar com outros pacotes gráficos como o ``pictex''
% Então é necessário usar apenas um dos pacotes conflitantes
\usepackage{tikz, tikz-cd}
\usepackage{epigraph}
\usetikzlibrary{babel}
\usetikzlibrary{positioning}
\usetikzlibrary{calc}
\usepackage[stix2]{newtxmath}
\usepackage{mathtools, nccmath}
%\usepackage[usenames,dvipsnames]{xcolor}
\usepackage{amssymb, mathrsfs}
\usepackage[numbers]{natbib}
\usepackage{fontspec}
\usepackage{enumitem} % pra substituir o símbolo do itemize por alguma coisa legal, tipo setinhas:  \begin{itemize}[label=\ding{212}]
\usepackage{pifont}
\usepackage{stmaryrd} % for \arrownot (notação de adjacência em grafo)
\usepackage{quiver} %pacote do quiver pra desenhar diagramas mais legais

%Fontes e Línguas
\usepackage[portuguese]{babel}
%\usepackage{eufrak}
\usepackage{emoji}
\setmainfont{DarwinSerif-Regular.otf}
\newfontfamily\fakebold[FakeBold=1.75]{DarwinSerif-Regular.otf}
\newfontfamily\fakeitalic[FakeSlant=0.2]{DarwinSerif-Regular.otf}
\renewcommand{\textbf}[1]{{\fakebold#1}}
\renewcommand{\textit}[1]{{\fakeitalic#1}}
\renewcommand{\bfseries}{\fakebold}
\renewcommand{\itshape}{\fakeitalic}



%Cores


%Comandinhos que gosto
\usepackage{mathtools}
\newcommand{\defeq}{\vcentcolon=}
\newcommand{\eqdef}{=\vcentcolon}
\newcommand{\bigslant}[2]{{\raisebox{.2em}{$#1$}\left/\raisebox{-.2em}{$#2$}\right.}}
\newcommand{\logo}{\Rightarrow}
\newcommand{\pulinho}{\vspace{.5cm}}
\newcommand{\edge}{%
  \mathrel{-}% minus as a relation
  \joinrel\joinrel % some backup
  \mathrel{-}% minus as a relation
}
\newcommand{\notedge}{%
  \mathrel{\mkern2mu}% a small advancement
  \arrownot % a short slash
  \mathrel{\mkern-2mu}% compensate
  \edge
}

\newcommand{\Aut}{\text{Aut}}
\newcommand{\Cay}{\text{Cay}}

\newcommand{\luafacil}{
\emoji{waxing-crescent-moon}
}

\newcommand{\luamedio}{
\emoji{first-quarter-moon}
}

\newcommand{\luadificil}{
\emoji{waxing-gibbous-moon}
}

\newcommand{\luaprojeto}{
\emoji{full-moon}
}

\let\varphi\phi

%Caixinhas coloridas

\theoremstyle{definition}
\newtheorem{theorem}{Teorema}
\newenvironment{teorema}{
\definecolor{shadecolor}{HTML}{f59aa3}
\begin{shaded}
\begin{theorem}
}
{
\end{theorem}
\end{shaded}
}

\theoremstyle{definition}
\newtheorem*{theorem*}{Teorema}
\newenvironment{teorema*}{
\definecolor{shadecolor}{HTML}{f59aa3}
\begin{shaded}
\begin{theorem*}
}
{
\end{theorem*}
\end{shaded}
}

\theoremstyle{definition}
\newtheorem{prop}{Proposição}
\newenvironment{proposicao}{
\definecolor{shadecolor}{HTML}{FFDBDF}
\begin{shaded}
\begin{prop}
}
{
\end{prop}
\end{shaded}
}

\theoremstyle{definition}
\newtheorem*{prop*}{Proposição}
\newenvironment{proposicao*}{
\definecolor{shadecolor}{HTML}{FFDBDF}
\begin{shaded}
\begin{prop*}
}
{
\end{prop*}
\end{shaded}
}

\theoremstyle{definition}
\newtheorem{cor}{Corolário}
\newenvironment{corolario}{
\definecolor{shadecolor}{HTML}{f59aa3}
\begin{shaded}
\begin{cor}
}
{
\end{cor}
\end{shaded}
}

\theoremstyle{definition}
\newtheorem*{cor*}{Corolário}
\newenvironment{corolario*}{
\definecolor{shadecolor}{HTML}{f59aa3}
\begin{shaded}
\begin{cor*}
}
{
\end{cor*}
\end{shaded}
}



\theoremstyle{definition}
\newtheorem{definition}{Definição}
\newenvironment{definicao}{
\definecolor{shadecolor}{HTML}{defcfc}
\begin{shaded}
\begin{definition}
}
{
\end{definition}
\end{shaded}
}

\theoremstyle{definition}
\newtheorem*{definition*}{Definição}
\newenvironment{definicao*}{
\definecolor{shadecolor}{HTML}{defcfc}
\begin{shaded}
\begin{definition*}
}
{
\end{definition*}
\end{shaded}
}

\theoremstyle{definition}
\newtheorem{example}{Exemplo}
\newenvironment{exemplo}{
\definecolor{shadecolor}{HTML}{b9e6d3}
\begin{shaded}
\begin{example}
}
{
\end{example}
\end{shaded}
}

\theoremstyle{definition}
\newtheorem*{example*}{Exemplo}
\newenvironment{exemplo*}{
\definecolor{shadecolor}{HTML}{b9e6d3}
\begin{shaded}
\begin{example*}
}
{
\end{example*}
\end{shaded}
}

\theoremstyle{definition}
\newtheorem{veronica}{Lema}
\newenvironment{lema}{
\definecolor{shadecolor}{HTML}{FFF1F2}
\begin{shaded}
\begin{veronica}
}
{
\end{veronica}
\end{shaded}
}

\theoremstyle{definition}
\newtheorem*{veronica*}{Lema}
\newenvironment{lema*}{
\definecolor{shadecolor}{HTML}{FFF1F2}
\begin{shaded}
\begin{veronica*}
}
{
\end{veronica*}
\end{shaded}
}

\newenvironment{prova}{\paragraph{Demonstração:}}{\hfill$\square$ 
\pulinho
}

\newenvironment{aviso}{
\definecolor{shadecolor}{HTML}{fffe9a}
\begin{shaded}
}{\end{shaded}}

\theoremstyle{definition}
\newtheorem{exersisio}{Exercício}
\newenvironment{exercicio}{
\definecolor{shadecolor}{HTML}{CFB5FF}
\begin{shaded}
\begin{exersisio}
}
{
\end{exersisio}
\end{shaded}
}

%Tamanho do papel
\usepackage[a4paper,top=1cm,bottom=1cm,margin=3.5cm]{geometry}



%%%%%%%%%%%%%%%%%%%%%%%%%%%%%%%%%%%%%%%%%%%%%%%


%Definições só pra esse texto
\usepackage[bb=boondox]{mathalpha}



\title{Quase livre}
\author{Violeta Mar}
\date{}


\begin{document}

\maketitle

Seja $G$ um grupo, e $X \subset G$ um subconjunto qualquer. Queremos construir um grupo que é `livre' sobre $X$ a menos de restrições existentes pela operação de $G$ em $X$.

Formalmente, vamos primeiro criar o grupo livre usual sobre $X$, chamemos de $F(X)$. Consideramos $X \subset F(X)$, ou seja, denotamos elementos de $X$ com o mesmo nome que seu correspondente em $F(X)$. Daí, definimos $Q_G \subset F(X)$ como sendo o subgrupo normal gerado por $\{xyz^-1 \in F(X) \mid x,y,z \in X \text{ e } xy = z\}$. Defina $F_G(X) \defeq \frac{F(X)}{Q_G}$, o grupo quase livre sobre $X$.

O grupo livre usual admite uma caracterização via uma propriedade universal. Pra obtermos uma propriedade universal, precisamos considerar funções que são `quase' homomorfismos, mais precisamente, funções que são homomorfismos sempre que a operação está bem definida em $X$.

\begin{definicao*}[Perfis]
  Seja $H$ algum outro grupo. Uma função $f: X \rightarrow H$ é dita um perfil se, sempre que $xy \in X$ para $x,y \in X$, temos que $f(xy) = f(x)f(y)$.
\end{definicao*}

\begin{proposicao*}[A propriedade universal da quase liberdade]
  Dado um perfil $f: X \rightarrow H$, existe um único homomorfismo $\tilde{f}: F_G(X) \rightarrow H$ que faz o seguinte diagrama comutar 
% https://q.uiver.app/#q=WzAsMyxbMCwxLCJYIl0sWzAsMCwiSCJdLFsxLDEsIkZfRyhYKSJdLFswLDIsImkiLDIseyJjb2xvdXIiOlsxNTEsNjIsNjBdfSxbMTUxLDYyLDYwLDFdXSxbMCwxLCJmIiwwLHsiY29sb3VyIjpbMTUxLDYyLDYwXX0sWzE1MSw2Miw2MCwxXV0sWzIsMSwiXFx0aWxkZXtmfSIsMix7ImNvbG91ciI6WzM1MCw0NCw2MF0sInN0eWxlIjp7ImJvZHkiOnsibmFtZSI6ImRhc2hlZCJ9fX0sWzM1MCw0NCw2MCwxXV1d
\[\begin{tikzcd}
	H & \\
	X & {F_G(X)}
	\arrow["f", color={rgb,255:red,90;green,216;blue,155}, from=2-1, to=1-1]
	\arrow["i"', color={rgb,255:red,90;green,216;blue,155}, from=2-1, to=2-2]
	\arrow["{\tilde{f}}"', color={rgb,255:red,198;green,108;blue,123}, dashed, from=2-2, to=1-1]
\end{tikzcd}\]
onde $i: X \rightarrow F_G(X)$ é a inclusão (denotamos com a cor verde funções que são perfis, e vermelho funções que são homomorfismos).
\end{proposicao*}
\begin{prova}
  A unicidade vem 
\end{prova}

Fazendo uma caça a diagrama usual, conseguimos provar que a propriedade universal caracteriza unicamente o grupo quase livre.

A última proposição também mostra que há uma bijeção entre perfis de $X$ em $H$ e homomorfismos entre $F_G(X)$ e $H$.


\end{document}
